% page 593 of the facsimile (image p-592 of the scan)
% block 180.3 x 246.3 mm ; left margin 21.7 mm
\pgc{0.000mm}{0.000mm}{
\l{\cel{57}585}
\l{}
\l{\cel{5}\sou{topazo}. (\sou{mineral.})\cel{2}Lapido diafana, orea briloze. - DEFIRS.}
\l{}
\l{\cel{5}\sou{topika}. (\sou{medic.}) Medikamento lokala, uzata extere (plastro,}
\l{\cel{8}kataplasmo, e c.) - DEFIRS.}
\l{}
\l{\cel{5}\sou{topinamburo}. (\sou{bot.}) Planto ek la familio \textquotedbl{}kompozaji\textquotedbl{}, di qua}
\l{\cel{8}la tuberkulo manjesas kom surogato di artichoko. L.\cel{1}\sou{halianthus}}
\l{\cel{8}\sou{tuberosus}. - DF.}
\l{}
\l{\cel{5}\sou{topografio}. La arto deskriptar e reprezentar mape la figuro}
\l{\cel{8}di loko, di la obstakli di tereno. - DEFIRS.}
\l{}
\l{\cel{5}\sou{topologio}. I. (\sou{geogr.)}\cel{2}Studio di la formi di la tereno, e di}
\l{\cel{8}la legi naturala (+leyi) qui guvernas li. - II. (\sou{geom.})}
\l{\cel{8}Fako di la geometrio da Riemann, en qua onu studias la ka-}
\l{\cel{8}rakteri di spaco - ta karakteri esante nur qualesala, ed omna}
\l{\cel{8}ideo pri mezuro eskartite. - DEFIS.}
\l{}
\l{\cel{5}\sou{toro.}\cel{2}I. (\sou{arkitekt.)}\cel{2}Muluro cilindra an la bazo di kolono, II.}
\l{\cel{8}(\sou{geom.)}\cel{2}Surfaco produktita per la rotaco di cirklo cirkum}
\l{\cel{8}rekto qua kontenesas en lua plano, ma qua ne pasas tra lua}
\l{\cel{8}centro. - DEFIS.}
\l{}
\l{\cel{5}\sou{torako}. (\sou{anat.})\cel{2}I. (\sou{che la vertebrozi}). Kavajo quan shirmas}
\l{\cel{8}parieto osta, e qua kontenas la organi precipua di la sango-}
\l{\cel{8}cirkulado e di la respirado. - II. (\sou{che la insekti)}. Parto di}
\l{\cel{8}la korpo inter la kapo e la abdomino, e sur qua insertesas}
\l{\cel{8}la ali e la gambi. - DEFIS.}
\l{}
\l{\cel{4}\sou{torcho}. Flamo-fonto grosiera, ek kordego tordita, indutita ye}
\l{\cel{8}rezino o vaxo - od ek bastono ek ligno rezinoza, kovrita ye}
\l{\cel{8}vaxo, sebo, quan onu portis lor la procesioni, la mariaji,}
\l{\cel{8}la funeri - quan onu tenis en sua manuo lor sua penitenco}
\l{\cel{8}publika o reparo-prego - e quan onu uzas ankore por lumizar}
\l{\cel{8}defilo dum nokto. - EFIS.}
\l{}
\l{\cel{5}\sou{tordar}. (\sou{trans.}) Kontraktar violentoze (korpo flexebla) sur}
\l{\cel{8}lu ipsa, per mantenar fixa un ek lua extremaji, en ica sinso,}
\l{\cel{8}e rotacigar la altra inversa-sinse. - EFIS.}
\l{}
\l{\cel{5}\sou{toreadoro}. Lor tauro-kombato, ta omna personi qui kombatas}
\l{\cel{8}kontre la tauro, en la areno : pikadoro, matadoro, e c.- DEFIRS}
\l{}
\l{\cel{5}\sou{torefaktar}. (\sou{trans.)}\cel{2}Submisar a la efiko da fairo intensa, qua}
\l{\cel{8}produktas komenco di karbonigo. - EFIS.}
\l{}
\l{\cel{5}\sou{torento}. Aquo-fluo, igita impetuoza, sive da la troa pento di la}
\l{\cel{8}tereno, sive da acenso efemera. - EFIS.}
\l{}
\l{\cel{5}\sou{torfo}. Materio kombustebla, sponjatra, nigratra, qua naskas}
\l{\cel{8}(maxim-multa-kaze) sub la aqui stagnanta, de la deskompozo}
\l{\cel{8}di la rezidui de vejetanti. - DEFIRS.}
\l{}
\l{\cel{5}\sou{torio}. (\sou{kemio}).Metalo skarsa, ek la grupo de la metali teroza,}
\l{\cel{8}quan Berzelius deskovris en 1928.\cel{1}\sou{Simb. kemiala}\cel{1}: Th. - DEFIRS}
}
